% Where the matrix links come from: group-valued links and their communicating parts.
% Drafted by Claude Opus 5.5 at the author's request, 2026-10-02. Numerical checks: link_origin_checks/.
\documentclass[11pt]{article}
\usepackage[margin=1in]{geometry}
\usepackage{amsmath,amssymb,amsthm}
\usepackage{booktabs}
\usepackage{array}
\usepackage[T1]{fontenc}
\usepackage{lmodern}
\usepackage{microtype}
\usepackage{xcolor}
\usepackage[colorlinks=true,linkcolor=blue!50!black,citecolor=green!40!black,urlcolor=blue!50!black]{hyperref}
\emergencystretch=2em

\newtheorem{proposition}{Proposition}
\newtheorem{lemma}[proposition]{Lemma}
\newtheorem{corollary}[proposition]{Corollary}
\theoremstyle{definition}
\newtheorem{definition}[proposition]{Definition}
\theoremstyle{remark}
\newtheorem{remark}[proposition]{Remark}

\newcommand{\GL}{\mathrm{GL}}
\newcommand{\gl}{\mathfrak{gl}}
\newcommand{\tr}{\operatorname{tr}}
\newcommand{\im}{\operatorname{im}}
\newcommand{\cfg}[1]{\texttt{#1}}

\title{Where the Matrix Links Come From:\\Group-Valued Links and Their Communicating Parts}
\author{Robert C. Dennis\\\small Technical note, drafted with Claude Opus 5.5, October 2, 2026}
\date{}

\begin{document}
\maketitle

\begin{abstract}
The projected attention routes of the gauge-covariant transformer communicate through matrix links $L_{ij}^{(h)}=U_iO_hT_hU_j^{-1}$, which have rank at most $d$ and are therefore not elements of $\GL(K)$. This note asks whether these links can be understood as parts of group-valued links $\Theta_{ij}\in\GL(K)$ that also contain a component that neither sends nor receives. A singular link cannot be a factor of a group element, so a literal product $\Theta=L\times N$ is impossible, but three exact constructions give the idea precise content. In the first, $L$ is a Lie-algebra element and $\Theta_{ij}=U_i\exp(X_{ij})U_j^{-1}$. The exponential of a rank-$d$ algebra element splits exactly as the flat transport plus a rank-$d$ channel term, $\Theta_{ij}=\Omega_{ij}+U_iO\varphi(TO)TU_j^{-1}$ with $\varphi(z)=(e^z-1)/z$, so the non-communicating part is the identity on the kernel of the sending map and $L$ is the first-order term. In the second, a primitive group link is compressed by frame-covariant projectors, and a block factorization relative to the parabolic subgroup that stabilizes the channel separates a communicating Levi block from a hidden Schur complement. In the third, the multi-head operator that the code already applies, $U_iOB_{ij}TU_j^{-1}$ with $B_{ij}$ the head-diagonal routing weights, is itself generically group-valued, and its loops measure a cocycle defect of the attention logits. Under Option A the routing belongs to per-head source labels outside the exponent, and requiring that only sent content cross between agents makes the receiver carry the complement. The flat term then becomes a residual on the unsent directions only; with tiled heads, as the code configures them, there is no residual and Option A reparameterizes the current head sum. The function $\varphi$ acts on weights rather than states, but the routing weights depend on the states, and a routed self-link yields a bounded block gate $e^{\beta_{ii}}$. For a channel of this form the linear link and the exponential link are two charts of the same embedded $\GL(d)$, related exactly by $O'=O\varphi(TO)$ wherever the principal logarithm exists, so the matrix link is the group-valued link in its affine chart and not a truncation of it; the charts differ at the boundary $\det=0$, in components, and in optimization geometry. We rank the three constructions by how principled they are and by how closely they match standard attention, and we record what each predicts and leaves open. All identities were checked numerically in float64 at small dimension; no model was trained.
\end{abstract}

\section{Setting and the question}
\label{sec:setting}

Agent $i$ carries a Gaussian belief $q_i=\mathcal N(\mu_i,\Sigma_i)$ on $\mathbb R^K$ and an invertible frame $U_i$, with reference moments $\xi_i=U_i^{-1}\mu_i$ and $C_i=U_i^{-1}\Sigma_iU_i^{-\top}$. A passive frame change $g_i\in\GL(K)$ acts by $(\mu_i,\Sigma_i,U_i)\mapsto(g_i\mu_i,g_i\Sigma_ig_i^\top,g_iU_i)$ and fixes $(\xi_i,C_i)$. The flat transport is $\Omega_{ij}=U_iU_j^{-1}$. Head $h$ of a layer owns a sending map $T_h:\mathbb R^K\to\mathbb R^d$ and a receiving map $O_h:\mathbb R^d\to\mathbb R^K$, and the code requires $K=Hd$. The matrix link is
\begin{equation}
L_{ij}^{(h)}=U_iO_hT_hU_j^{-1},\qquad L_{ij}^{(h)}\mapsto g_iL_{ij}^{(h)}g_j^{-1},
\label{eq:link}
\end{equation}
and the projected predictive mean is $\mu_i^+=r\mu_i+\sum_{h,j}\beta_{ij}^{(h)}L_{ij}^{(h)}\mu_j$ with $r\in\{0,1\}$ set by the attention residual \cite{dennis2026projpred}. Stacking the heads gives square matrices $O=[O_1\cdots O_H]$ and $T=[T_1;\dots;T_H]$, both $K\times K$, and the head-diagonal routing matrix $B_{ij}=\bigoplus_h\beta_{ij}^{(h)}I_d$.

A group-valued link is a field $\Theta_{ij}\in\GL(K)$ with the same law $\Theta_{ij}\mapsto g_i\Theta_{ij}g_j^{-1}$, which is the lattice gauge field of Wilson \cite{wilson1974}. The question is whether $L_{ij}$ can be read as the part of some $\Theta_{ij}$ that takes part in sending and receiving, with a remainder that does not. To keep notation apart from the sending maps $T_h$, we write $N$ for that remainder.

Two facts constrain any answer. First, $L_{ij}^{(h)}$ has rank at most $d<K$, so for any $M$ the product $L^{(h)}M$ is singular, and no factorization $\Theta=L^{(h)}N$ of a group element exists. The relation between $L$ and $\Theta$ must be additive, through a block of $\Theta$, or infinitesimal, through the Lie algebra. Second, on causal support every source satisfies $j\le i$, so a closed directed walk can only be a self-loop, a point raised in the review of the attention manuscript \cite{glkreview2026}. Loop observables with more than one edge need reversed edges, which in lattice gauge theory are traversed with the inverse link \cite{wilson1974,kogut1979}. Inverses exist for group-valued links and not for singular ones, so a group-valued origin for $L$ is also what gives the link field nontrivial loop observables on causal support.

\section{Option A: the link is a Lie-algebra element}
\label{sec:exp}

\begin{definition}[Exponential link]
For an algebra field $X_{ij}\in\gl(K)$ written in reference coordinates, set
\begin{equation}
\Theta_{ij}=U_i\exp(X_{ij})U_j^{-1}.
\label{eq:exp_link}
\end{equation}
\end{definition}

The law $\Theta_{ij}\mapsto g_i\Theta_{ij}g_j^{-1}$ holds for every $X_{ij}$, because $X_{ij}$ is expressed in reference coordinates and the frames carry all of the passive action. The identity $\det\exp(X)=e^{\tr X}$ makes $\Theta_{ij}$ invertible for every algebra element, singular ones included. The choice that reproduces the current links is the rank-$d$ algebra element $X=O_hT_h$ for one head, or $X=\sum_hO_hT_h=OT$ for all heads with equal weight.

\begin{lemma}[Exponential of a low-rank element]
\label{lem:lowrank_exp}
Let $P$ be $K\times r$ and $Q$ be $r\times K$, and let $\varphi(Z)=\sum_{n\ge1}Z^{n-1}/n!$, which equals $(e^Z-I)Z^{-1}$ when $Z$ is invertible. Then
\begin{equation}
\exp(PQ)=I_K+P\varphi(QP)Q.
\label{eq:lowrank_exp}
\end{equation}
\end{lemma}

\begin{proof}
For $n\ge1$, $(PQ)^n=P(QP)^{n-1}Q$ by regrouping the product. Summing the exponential series term by term gives $I_K+P\big(\sum_{n\ge1}(QP)^{n-1}/n!\big)Q$, and the inner series converges for every $r\times r$ matrix.
\end{proof}

\begin{proposition}[Flat transport plus a channel term]
\label{prop:split}
With $X_{ij}=O_hT_h$,
\begin{equation}
\Theta_{ij}=\Omega_{ij}+\widetilde L_{ij}^{(h)},\qquad
\widetilde L_{ij}^{(h)}=U_iO_h\varphi(T_hO_h)T_hU_j^{-1}.
\label{eq:split}
\end{equation}
The following hold.
\begin{enumerate}
\item $\exp(O_hT_h)$ acts as the identity on $\ker T_h$, so on the reference subspace that the head does not read, $\Theta_{ij}$ coincides with the flat transport $\Omega_{ij}$.
\item The image of $\exp(O_hT_h)-I$ lies in $\im O_h$, so $\Theta_{ij}$ differs from $\Omega_{ij}$ only in the directions that the head writes.
\item Under a scaling $X_{ij}=sO_hT_h$, $\Theta_{ij}=\Omega_{ij}+sL_{ij}^{(h)}+O(s^2)$, so the matrix link of Eq.~\eqref{eq:link} is the first-order term.
\item $\det\Theta_{ij}=\det(U_iU_j^{-1})e^{\tr(T_hO_h)}$.
\end{enumerate}
\end{proposition}

\begin{proof}
Equation~\eqref{eq:split} is Lemma~\ref{lem:lowrank_exp} with $P=O_h$ and $Q=T_h$, conjugated by the frames. Items 1 and 2 are read off Eq.~\eqref{eq:lowrank_exp}: the correction is $O_h(\cdots)T_h$. Item 3 follows from $\varphi(sZ)=I+O(s)$. Item 4 is $\det\exp(X)=e^{\tr X}$ with $\tr(O_hT_h)=\tr(T_hO_h)$.
\end{proof}

Proposition~\ref{prop:split} gives the original question an exact answer. The remainder $N$ of the question is the identity on the non-communicating subspace $\ker T_h$, and the communicating part is $\widetilde L_{ij}^{(h)}$. That part equals the matrix link to first order and carries the corrections $\varphi(T_hO_h)$ at higher order. Nothing in $\Theta_{ij}$ is hidden: every entry is fixed by the frames and by $(O_h,T_h)$.

\paragraph{The two historical routes as orders of one link.}
Contracting Eq.~\eqref{eq:split} against routing weights gives the message
\begin{equation}
\sum_j\beta_{ij}\Theta_{ij}\mu_j
=\underbrace{\sum_j\beta_{ij}\Omega_{ij}\mu_j}_{\text{flat transport}}
+\underbrace{\sum_j\beta_{ij}L_{ij}\mu_j}_{\text{projected link}}+O(\|T_hO_h\|).
\label{eq:two_routes}
\end{equation}
The first term is the transported message of the earlier transported route, and the second is the projected predictive message. Option A therefore places both routes inside one group-valued link, as its zeroth- and first-order terms. This is an observation about the form of the message, not a claim that either retired objective is recovered.

\paragraph{Routing inside the exponent.}
Setting $X_{ij}=OB_{ij}T=\sum_h\beta_{ij}^{(h)}O_hT_h$ instead dresses the algebra element with the routing weights, and Lemma~\ref{lem:lowrank_exp} with $P=O$ and $Q=B_{ij}T$ gives
\begin{equation}
\exp(OB_{ij}T)=I+O\varphi(B_{ij}TO)B_{ij}T.
\label{eq:exp_routed}
\end{equation}
The first-order term is the multi-head operator of Section~\ref{sec:headsum}. Heads interact at second order through the cross-head blocks of $TO$. With routing inside the exponent, the weights no longer normalize a sum over sources, so the aggregation over $j$ must be chosen separately. Section~\ref{sec:aggregation} resolves this by keeping the routing outside the exponent.

\paragraph{Loops.}
For a triangle with a reversed edge, the loop is $\Theta_{ij}\Theta_{jk}\Theta_{ik}^{-1}=U_ie^{X_{ij}}e^{X_{jk}}e^{-X_{ik}}U_i^{-1}$, and its conjugacy class, for example its trace, is frame invariant. When the algebra elements commute, which happens for $X=OB T$ with $TO=I$ since then $X_{ij}=OB_{ij}O^{-1}$, the loop equals
\begin{equation}
U_iO\exp\big(B_{ij}+B_{jk}-B_{ik}\big)O^{-1}U_i^{-1},
\label{eq:exp_loop}
\end{equation}
so the curvature is the per-head additive defect $\beta_{ij}^{(h)}+\beta_{jk}^{(h)}-\beta_{ik}^{(h)}$. In general the loop follows the Baker-Campbell-Hausdorff series, and the commutators $[X_{ij},X_{jk}]$ make it nonabelian.

\subsection{Linear and exponential links are two charts of one group}
\label{sec:charts}

The phrase ``first-order term'' in Proposition~\ref{prop:split} invites the reading that the matrix link is a weak-field truncation of the group link. For a channel of this form the truncation loses nothing. Fix $O\in\mathbb R^{K\times d}$ and $T\in\mathbb R^{d\times K}$ with $TO$ invertible, so that $\mathbb R^K=W\oplus W'$ with $W=\im O$ and $W'=\ker T$, and let $P=O(TO)^{-1}T$ be the projector onto $W$ along $W'$. For $B\in\GL(d)$ define
\begin{equation}
\iota(B)=(I-P)+O\,B\,(TO)^{-1}T,
\label{eq:embed}
\end{equation}
the operator that fixes $W'$ pointwise and acts on $W$ as $B$ in the basis formed by the columns of $O$.

\begin{proposition}[Two charts of one group]
\label{prop:charts}
\begin{enumerate}
\item $\iota:\GL(d)\to\GL(K)$ is an injective homomorphism with $\det\iota(B)=\det B$. Its image $\mathcal G_W$ is the Levi factor of the stabilizer of the splitting $W\oplus W'$ that acts trivially on $W'$.
\item \emph{Linear links.} For $C\in\mathbb R^{d\times d}$, $I+OCT=\iota(I_d+C\,TO)$, so $\det(I+OCT)=\det(I_d+C\,TO)$.
\item \emph{Exponential links.} $\exp(OT)=\iota(e^{TO})$.
\item \emph{Exact conversion.} Let $O'\in\mathbb R^{K\times d}$ be such that $I_d+TO'$ has no eigenvalue on the closed negative real axis $(-\infty,0]$. Put $Z=\log(I_d+TO')$ (principal logarithm) and $O=O'\varphi(Z)^{-1}$. Then $\varphi(Z)$ is invertible, $TO=Z$, and
\begin{equation}
\exp(OT)=I_K+O'T.
\label{eq:conversion}
\end{equation}
\end{enumerate}
\end{proposition}

\begin{proof}
For item 1, $\iota(B)Oy=OBy$ because $(I-P)O=0$ and $(TO)^{-1}TO=I$, and $\iota(B)w'=w'$ for $w'\in W'$ because $Pw'=0$. An operator on $W\oplus W'$ is determined by its action on the two summands, so $\iota(B)\iota(B')=\iota(BB')$ and $\det\iota(B)=\det B\cdot1$. For item 2, $(I+OCT)Oy=O(I+C\,TO)y$ and $(I+OCT)w'=w'$. For item 3, $(OT)^nO=O(TO)^n$ gives $\exp(OT)Oy=Oe^{TO}y$, and $\exp(OT)w'=w'$. For item 4, the eigenvalues of $Z$ have imaginary part in $(-\pi,\pi)$, so $\varphi(Z)$, whose eigenvalues $(e^z-1)/z$ vanish only at $z\in2\pi i\mathbb Z\setminus\{0\}$, is invertible. Since $e^Z-I=TO'=Z\varphi(Z)$, we get $TO=TO'\varphi(Z)^{-1}=Z$, and Lemma~\ref{lem:lowrank_exp} gives $\exp(OT)=I+O\varphi(TO)T=I+O'\varphi(Z)^{-1}\varphi(Z)T=I+O'T$.
\end{proof}

Item 4 says that the flat transport plus a linear channel term, $\Omega_{ij}+U_iO'TU_j^{-1}$, is exactly an exponential link for every $O'$ in an open set that contains $O'=0$. The two families are the same subgroup written in an affine chart ($B=I+C\,TO$) and a logarithmic chart ($B=e^{Z}$). The expansion $\Theta=\Omega+L+O(X^2)$ of Proposition~\ref{prop:split} is the statement that the charts agree to first order at the identity, and with $\widetilde L=O\varphi(TO)T=\exp(OT)-I$ in reference coordinates the series $X=\log(I+\widetilde L)=\widetilde L-\tfrac12\widetilde L^2+\cdots$ converts one chart to the other. It is not a statement that the linear link approximates the group link. The same holds for the lattice relation $U=e^{aA}$, where the algebra-valued potential $A$ is the logarithmic coordinate, with the difference that here the frames dress the potential, so that $X$ is frame invariant in reference coordinates.

\begin{corollary}[Tiled heads]
\label{cor:tiled_chart}
If the stacked $O$ and $T$ are square and invertible, so that $W=\mathbb R^K$, the operator $M=O'T$ that the projected predictive route applies at fixed routing is itself a group element, $M\in\GL(K)$ when $\det M\ne0$, and $M=\exp(X)$ for $X=OT$ exactly when $M$ has a real logarithm, with $O=XT^{-1}$. The code's identity initialization $M=I$ is the origin $X=0$ of the logarithmic chart. The product of $r$ links along a path with a shared generator is $M^r=e^{rX}$, a one-parameter subgroup.
\end{corollary}

The corollary uses $M^r=e^{rX}$ to restate the path-composition remark of the matrix-link construction: the $r$-edge path $U_{i_0}M^rU_{i_r}^{-1}$ is the group element of a one-parameter subgroup, not a separate object.

\paragraph{What the choice of chart does change.}
First, the affine chart is the whole matrix space. A step in $(O',T)$ can carry $\det(I_K+O'T)=\det(I_d+TO')$ through zero (in the setting of Corollary~\ref{cor:tiled_chart}, $\det M$), where the link becomes singular and the information carried by the $W$ component is destroyed. The logarithmic chart cannot reach that boundary, because $\det e^Z=e^{\tr Z}>0$. The same argument shows that the exponential chart covers only matrices with a real logarithm: it excludes negative-determinant components, which an optimizer in the affine chart reaches only by passing through $\det=0$, and a real invertible matrix has a real logarithm only if its Jordan blocks for negative eigenvalues occur in pairs \cite{culver1966}. Second, the two charts give different optimization geometries, since the code's M-step takes Euclidean-type steps on $(O',T)$, which is the affine chart. Third, they differ at order one for finite generators: at $X=I$, $e^X=eI$ while $I+X=2I$. Fourth, these statements concern an unrouted operator with a single generator. The routed head sum $\sum_h\beta^{(h)}O_hT_h$ weights heads differently and does not lie in the one-generator family, and the hidden block and edge-dependent generators of Option B remain outside $\mathcal G_W$.

\section{Option B: the link is a compression of a primitive group link}
\label{sec:compression}

Here the group-valued link is taken as primary, $\Theta_{ij}=U_iG_{ij}U_j^{-1}$ with $G_{ij}\in\GL(K)$, and the channel maps act as covariant projectors. When $T_hO_h$ is invertible, define the reference projector $P_h=O_h(T_hO_h)^{-1}T_h$, which projects onto $\im O_h$ along $\ker T_h$, and its frame-dressed version $\Pi_{i,h}=U_iP_hU_i^{-1}$, which transforms by conjugation. Every link then decomposes into four covariant blocks,
\begin{equation}
\Theta_{ij}=\Pi_{i}\Theta_{ij}\Pi_{j}+\Pi_{i}\Theta_{ij}(I-\Pi_{j})+(I-\Pi_{i})\Theta_{ij}\Pi_{j}+(I-\Pi_{i})\Theta_{ij}(I-\Pi_{j}),
\label{eq:blocks}
\end{equation}
where the first block is the communicating link, the last is the non-communicating remainder $N$ of the question, and the middle two mix them. More generally, distinct receive and send projectors give the oblique compression $\Pi^R_i\Theta_{ij}\Pi^S_j$, and any rank-$d$ map whose range lies in $\im O_h$ and whose kernel contains $\ker T_h$ arises this way from some invertible $G_{ij}$.

\begin{proposition}[Block factorization]
\label{prop:levi}
In a reference basis adapted to $\mathbb R^K=W\oplus W'$ with $W=\im O_h$, write $G=\begin{pmatrix}A&B\\C&D\end{pmatrix}$. If $A$ is invertible, then
\begin{equation}
G=\begin{pmatrix}I&0\\CA^{-1}&I\end{pmatrix}
\begin{pmatrix}A&0\\0&S\end{pmatrix}
\begin{pmatrix}I&A^{-1}B\\0&I\end{pmatrix},\qquad S=D-CA^{-1}B,
\label{eq:levi}
\end{equation}
and $\det G=\det A\det S$.
\end{proposition}

The middle factor lies in $\GL(d)\times\GL(K-d)$, the Levi factor of the parabolic subgroup that stabilizes $W$ \cite{humphreys1975}. Its first block $A$ is the communicating link, and $S$, the Schur complement \cite{zhang2005schur}, is the hidden part with the communicating directions eliminated. The outer unipotent factors carry the mixing. This is the precise multiplicative form of $\Theta=L\times N$: the product holds between the outer mixing factors and the block-diagonal pair $(A,S)$, not between $L$ and $N$ directly.

\begin{remark}[Reduction of structure group]
Assigning the subspace $W_i=U_i\im O_h$ to every agent is a section of the associated Grassmannian bundle, and it reduces the structure group from $\GL(K)$ to the parabolic stabilizer. Reductions of this kind are how a Higgs field breaks a gauge symmetry. The analogy is formal; no potential or dynamics for $W_i$ is proposed here.
\end{remark}

\paragraph{What Option B needs and what it cannot see.}
The current model is the case $G_{ij}=G$ constant, with every head reading its own block. A genuine extension needs a model for $G_{ij}$: a learned relative-position element $G_{i-j}$, a state-dependent element $G(\xi_i,\xi_j)$, or a dynamical link with its own energy. Rotary embeddings are the closest precedent in standard practice, a relative-position orthogonal link inserted between queries and keys \cite{su2021roformer}. The code's \cfg{rope\_on\_value=True} also dresses the value link with the rotation, $V_iO_hT_hV_j^{-1}=U_iR_i^\top O_hT_hR_jU_j^{-1}$. The decomposition also says what cannot be learned. The blocks $B$, $C$ and $D$ never enter a message, so they receive no gradient from the routing or content path unless another term reads them. The projected gradient analysis gives one such term \cite{dennis2026projgrad}. The ambient-alignment receiver energy pins the complement of each output range to zero, while the channel-likelihood energy leaves that complement at the prior. Under Option B, the identifiability of the hidden block is therefore a property of the chosen receiver energy, not of the link.

\section{Option C: the multi-head operator is already group-valued}
\label{sec:headsum}

\begin{proposition}[Head-sum link]
\label{prop:headsum}
The edge operator that the projected routes apply at fixed routing,
\begin{equation}
\mathcal C_{ij}=\sum_h\beta_{ij}^{(h)}L_{ij}^{(h)}=U_iOB_{ij}TU_j^{-1},
\label{eq:headsum}
\end{equation}
transforms as a link and lies in $\GL(K)$ whenever the stacked $O$ and $T$ are invertible and every $\beta_{ij}^{(h)}>0$. Its determinant is $\det(U_iU_j^{-1})\det(TO)\prod_h(\beta_{ij}^{(h)})^d$.
\end{proposition}

\begin{proof}
The identity $\sum_h\beta^{(h)}O_hT_h=OB T$ is block multiplication with $B=\bigoplus_h\beta^{(h)}I_d$, and the determinant is multiplicative.
\end{proof}

Invertibility of $O$ and $T$ is generic and holds at the code's identity initialization, where $O_h$ and $T_h^\top$ are the coordinate block embeddings. Softmax weights are strictly positive. Proposition~\ref{prop:headsum} is the remark in the attention manuscript that the multi-head operator can be full rank \cite{dennis2026glkv2}, made into a statement about group values. The heads are then the decomposition the question asks for: each head-diagonal block of $\mathcal C_{ij}$ is one head's link, the routing weights are the diagonal factor $B_{ij}$, and $TO$ carries the mixing.

This is also exactly the value path of standard multi-head attention. With trivial frames, the multi-head output of a transformer is $\sum_hW_O^{(h)}\sum_j\beta_{ij}^{(h)}W_V^{(h)}x_j=\sum_jW_OB_{ij}W_Vx_j$ \cite{vaswani2017}, the frame-free case of Eq.~\eqref{eq:headsum} with $O=W_O$ and $T=W_V$.

\paragraph{Loops measure a cocycle defect of the logits.}
With $TO=I$ the head blocks commute, and the triangle loop with one reversed edge is
\begin{equation}
\mathcal C_{ij}\mathcal C_{jk}\mathcal C_{ik}^{-1}
=U_iO\Big(\bigoplus_h\frac{\beta_{ij}^{(h)}\beta_{jk}^{(h)}}{\beta_{ik}^{(h)}}I_d\Big)O^{-1}U_i^{-1}.
\label{eq:headsum_loop}
\end{equation}
For softmax routing $\beta_{ij}=e^{s_{ij}}/Z_i$, the logarithm of each head factor is
\begin{equation}
\log\frac{\beta_{ij}\beta_{jk}}{\beta_{ik}}=s_{ij}+s_{jk}-s_{ik}-\log Z_j,
\label{eq:cocycle}
\end{equation}
because the receiver normalizer $Z_i$ cancels between $\beta_{ij}$ and $\beta_{ik}$. The loop is trivial exactly when the logits form a cocycle up to the intermediate normalizer. In the projected routes the routing is $\beta_{ij}^{(h)}=\pi_{ij}^{(h)}\exp(-D_{ij}^{(h)}/\kappa_h)/Z_i^{(h)}$, with $D_{ij}^{(h)}$ the Gaussian KL between head $h$'s query law of $i$ and key law of $j$, so $s_{ij}=\log\pi_{ij}-D_{ij}/\kappa$ and
\begin{equation}
\log\frac{\beta_{ij}\beta_{jk}}{\beta_{ik}}=-\frac{D_{ij}+D_{jk}-D_{ik}}{\kappa}+\log\frac{\pi_{ij}\pi_{jk}}{\pi_{ik}}-\log Z_j.
\label{eq:kl_triangle}
\end{equation}
The leading term is the triangle defect of the routing costs. A metric would make $D_{ij}+D_{jk}-D_{ik}$ nonnegative, but the Kullback-Leibler divergence obeys no triangle inequality, and here it compares a query law with a key law under different maps $S_Q$ and $S_K$, so the defect can take either sign. With uniform priors, Eq.~\eqref{eq:kl_triangle} says that the loop of the head-sum link measures how far the routing geometry is from additive along the path $i\to j\to k$, offset by the intermediate token's log normalizer. For dot-product logits $s_{ij}=q_i^\top k_j/\sqrt d$ the defect is instead a bilinear expression in the queries and keys of the three tokens. Reading attention weights as an abelian connection is an interpretation suggested by Eq.~\eqref{eq:headsum_loop}; Proposition~\ref{prop:headsum} establishes only the algebra.

\paragraph{Limits.}
Group-valuedness in Option C is generic rather than structural. The condition number of $\mathcal C_{ij}$ grows with the ratio of the largest to the smallest head weight on the edge, so concentrated routing pushes it toward the boundary of $\GL(K)$. A singular $O$ or $T$, or $K\ne Hd$, removes group-valuedness altogether.

\section{Aggregation, the residual, and the nonlinearity under Option A}
\label{sec:aggregation}

This section settles how sources combine under Option A, shows when a residual follows from the construction and when it does not, and asks whether the function $\varphi$ is related to the activations of a transformer MLP.

\subsection{The complement split}

Stack the heads of a layer into $O$ of size $K\times Hd$ and $T$ of size $Hd\times K$, allow $Hd\le K$, and assume $TO$ is invertible. Let $P=O(TO)^{-1}T$, the projector onto $W=\im O$ along $W'=\ker T$.

\begin{lemma}[Complement split]
\label{lem:complement}
\begin{equation}
\exp(OT)=(I-P)+O\,E\,T,\qquad E=e^{TO}(TO)^{-1}.
\label{eq:complement}
\end{equation}
\end{lemma}

\begin{proof}
By Lemma~\ref{lem:lowrank_exp}, $\exp(OT)=I+O\varphi(TO)T=(I-P)+O\big[(TO)^{-1}+\varphi(TO)\big]T$, and $(TO)^{-1}+\varphi(TO)=(TO)^{-1}+(e^{TO}-I)(TO)^{-1}=e^{TO}(TO)^{-1}$ because every term is a function of $TO$.
\end{proof}

The first term is the non-communicating part of the question: it is the identity on $W'$ and annihilates $W$. The second term reads the $Hd$ sent coordinates $T\xi$, mixes them through the $Hd\times Hd$ matrix $E$, and writes them into $W$. Equation~\eqref{eq:complement} is exact, so the group-valued link separates into a complement transport and a channel transport with no remainder.

\subsection{Routing belongs outside the exponent}

A link acts on one source state. A mixture over sources therefore needs the routing weights as probabilities over which state each part of the link carries, not as coefficients inside the generator. With $X_{ij}=OB_{ij}T$ in the exponent, the routed links $e^{X_{ij}}$ for different $j$ act on different vectors and no single group element aggregates them: a product $\prod_je^{X_{ij}}$ has no source vector to act on, and a sum of exponentials is not normalized by $\sum_j\beta_{ij}=1$. We therefore take the generator to be routing independent, $X=OT$ in reference coordinates, and attach the routing to source labels. The routed generator keeps its use as a loop observable (Eq.~\eqref{eq:exp_loop}) but does not define a message.

\subsection{The aggregation rule}

Following the predictive kernel of the projected predictive route \cite{dennis2026projpred}, each head draws its own source label $J^h$ with probability $\beta_{ij}^{(h)}$, independently across heads. In Eq.~\eqref{eq:complement} each sent block $T_h\xi$ is read from its head's chosen source, and $E$ mixes the sent blocks into $W$. The complement term needs a carrier of its own, a label $J^0$. The transported reference state is
\begin{equation}
\xi_i'=(I-P)\xi_{J^0}+\sum_hO\,E_{:,h}\,T_h\xi_{J^h},
\label{eq:transported_state}
\end{equation}
where $E_{:,h}$ is the block column of $E$ that receives head $h$. The expectation over the independent labels is
\begin{equation}
\mathbb E[\xi_i']=\mathbb E[(I-P)\xi_{J^0}]+\sum_h\sum_j\beta_{ij}^{(h)}O\,E_{:,h}\,T_h\xi_j.
\label{eq:aggregate}
\end{equation}
Each sent block travels with its own head's label, so the cross-head mixing in $E$ does not couple the head routings. Only the complement carrier remains to be chosen. The weights are functions of the states, $\beta_{ij}^{(h)}=\beta_{ij}^{(h)}(\mu,\Sigma)$, through the Gaussian comparison costs of queries and keys. Both projected routes evaluate them at the layer's input states and hold them fixed while the message is formed, which is exact inference in the declared cut model of the projected gradient analysis \cite{dennis2026projgrad}. Equation~\eqref{eq:aggregate} is therefore the expectation conditional on those states, and the state dependence of $\beta$ does not alter the aggregation rule.

\begin{proposition}[The communication principle fixes the complement carrier]
\label{prop:self_complement}
If only sent content may cross between agents, the complement must be carried by the receiver, $J^0=i$, and
\begin{equation}
\mu_i^+=U_i(I-P)U_i^{-1}\mu_i+\sum_h\sum_j\beta_{ij}^{(h)}U_iO\,E_{:,h}\,T_hU_j^{-1}\mu_j.
\label{eq:self_complement}
\end{equation}
\end{proposition}

\begin{proof}
The complement $W'=\ker T$ holds exactly the directions that no head sends. A carrier $J^0=j\ne i$ would place agent $j$'s unsent coordinates in agent $i$'s state, which contradicts the premise, so the receiver is the only admissible carrier. Equation~\eqref{eq:aggregate} with $J^0=i$, conjugated by the frames, gives the formula.
\end{proof}

The alternative, a complement carried by a routed label $J^0\sim\beta^{(0)}$, gives $\sum_j\beta_{ij}^{(0)}(I-P)\xi_j$, which moves the unsent coordinates of other agents. That is a flat-transport message of the kind the earlier transported route used, and it is admissible only if agents are meant to read each other's full state.

\subsection{Does the flat term produce a residual?}

Proposition~\ref{prop:self_complement} answers the question. Under the communication principle the flat term becomes a residual, but only on the complement: the receiver keeps its own state exactly in the directions that no head writes and takes the transported message in the directions that heads do write. This projected residual $U_i(I-P)U_i^{-1}\mu_i$ differs from the full residual $r\mu_i$ that the code adds under \cfg{attention\_residual}.

\begin{corollary}[Tiled heads leave no residual]
\label{cor:tiled}
If the heads tile the fiber, $Hd=K$ with $O$ and $T$ invertible, then $P=I$, the complement is empty, and Eq.~\eqref{eq:self_complement} reduces to
\begin{equation}
\mu_i^+=\sum_h\sum_j\beta_{ij}^{(h)}U_iO_h'T_hU_j^{-1}\mu_j,\qquad O_h'=O\,E_{:,h},
\label{eq:tiled}
\end{equation}
the residual-free projected predictive head sum with receiving maps $O_h'$.
\end{corollary}

The code requires $K=Hd$, so with invertible stacked maps its heads tile the fiber, and Option A then implies no residual at all. Two consequences follow. First, the $r=1$ residual of the code has no counterpart in Option A under the configuration the code uses, and the receiver's own state reaches $\mu_i^+$ only through the self-weights $\beta_{ii}^{(h)}$. Second, in the tiled case Option A adds no expressive power to the predictive route: its head sum lies in the current function class, with receiving maps of the special form $O\,E_{:,h}$. We do not claim the converse, that every current set of receiving maps arises this way. Until the generator is made edge dependent, the content of Option A in the tiled case is the interpretation together with the loop observables of the group-valued link.

A residual of the projected kind requires channels that do not tile, $Hd<K$. Equation~\eqref{eq:self_complement} then describes an architecture in which $K-Hd$ coordinates form a residual stream that no head writes and every head can read through $T$. In this reading the residual stream is the complement of the communication channels.

\paragraph{The same answer from Gaussian fusion.}
The projected gradient route reaches the same structure by a different argument \cite{dennis2026projgrad}. Its channel-likelihood receiver keeps the prior in directions that no head observes, and in the tiled limit with a flat prior and equal receiving noise it recovers the residual-free predictive head sum. At finite precision, fusion also yields a partial residual in the observed directions. For an isotropic prior of precision $\alpha$ and message evidence of precision $\lambda$ on $W$, the posterior mean is
\begin{equation}
\xi_i^+=(I-P_W)\xi_i+P_W\frac{\alpha\xi_i+\lambda m_i}{\alpha+\lambda},
\label{eq:fusion_residual}
\end{equation}
with $P_W$ the orthogonal projector onto $W$ and $m_i$ the message. The residual coefficient is the prior's share of precision, $\alpha/(\alpha+\lambda)$, in the observed directions and one in the unobserved directions, and Eq.~\eqref{eq:self_complement} is the limit $\lambda/\alpha\to\infty$. Both derivations give a residual that is full on unobserved directions and partial or zero on observed ones. The additive residual of the code, $\mu_i+m_i$, is neither: it adds the full own state to the message in every direction without normalization. That form is the residual connection of standard transformers \cite{he2016resnet}, whose justification is optimization rather than inference.

\subsection{Is \texorpdfstring{$\varphi$}{phi} an activation function?}

Not with the routing-independent generator $X=OT$. There $\varphi$ acts on the $d\times d$ round-trip matrix $T_hO_h$, which is built from weights, the link is linear in the transported state, and in the tiled case $\varphi$ is absorbed into the receiving maps (Corollary~\ref{cor:tiled}). The message is still nonlinear in the states, but only through $\beta(\mu,\Sigma)$, the way standard attention is nonlinear through its softmax.

The answer changes once the generator carries the routing. Because $\beta_{ij}^{(h)}=\beta_{ij}^{(h)}(\mu,\Sigma)$, the routed generator $X_{ij}=OB_{ij}T$ depends on the states, and $\varphi(B_{ij}TO)$ in Eq.~\eqref{eq:exp_routed} becomes a matrix function of them. On the self-link this is a well-defined single-source transformation, so the objection of the previous subsection to routing inside the exponent does not apply there. For tiled heads with $TO=I$,
\begin{equation}
\exp(OB_{ii}T)=\sum_he^{\beta_{ii}^{(h)}(\mu,\Sigma)}O_hT_h,
\label{eq:self_gate}
\end{equation}
so each head's block of the receiver state is multiplied by $e^{\beta_{ii}^{(h)}}\in(1,e)$, a gate whose argument is the self-routing weight $\beta_{ii}^{(h)}=\pi_{ii}^{(h)}e^{-D_{ii}^{(h)}/\kappa_h}/Z_i^{(h)}$. Here $D_{ii}$ compares the token's own query and key laws, which differ because $S_Q\ne S_K$, and $Z_i$ compares that cost with every competing source. This is a state-dependent nonlinearity produced by the group structure itself. It differs from a GELU in three ways. The gate reads the state only through the self-routing weight, which compares the token with its sources rather than reading the token's own coordinates. It acts on whole head blocks rather than coordinatewise. And it lies between $1$ and $e$, so it can amplify a block by at most a factor of $e$ but can never suppress it or change its sign.

A general state-dependent generator on the self-link, $\Theta_{ii}=U_i\exp(X(\xi_i))U_i^{-1}$, does produce the gated form of modern activations. For a diagonal generator $X(\xi)=\mathrm{diag}(a(\xi))$ the update is $\xi_k\mapsto\xi_ke^{a_k(\xi)}$, a coordinatewise gate. Every positive gate has this form, since $x\sigma(x)=\exp(\log\sigma(x))x$, so the SiLU $x\sigma(x)$ \cite{elfwing2018silu} and the GELU $x\Phi(x)$ \cite{hendrycks2016gelu} are exponential self-links with generators $\log\sigma$ and $\log\Phi$. In lattice language the attention links are edge variables and a gated MLP is a state-dependent site term.

This is a rewriting, not a derivation. It places the MLP in the same group-valued language as the links but does not explain why the generator should be $\log\sigma$ or $\log\Phi$. A derivation needs a principle that fixes the generator. Template routing supplies one for a normalized gate, replacing independent two-way competitions by one softmax over Gaussian templates derived from a mixture prior \cite{dennis2026template}.

\section{How the three options relate}
\label{sec:relations}

The three constructions nest. Option C is Option B with the reference element $G_{ij}=OB_{ij}T$ and the head projectors. Option A with routing inside the exponent agrees with Option C at first order, $\exp(OB_{ij}T)=I+OB_{ij}T+O(\|B_{ij}\|^2)$, apart from the identity, which contributes the flat transport. Option B is the most general frame: it holds any group link together with any choice of channel projectors. Option A restricts the link to the exponential of the channel algebra element and so fixes the hidden block to the identity. Option C is the model as it stands, read as a link.

\section{Ranking}
\label{sec:ranking}

Table~\ref{tab:principled} orders the options by how principled they are. Here the criteria are that the group property should be structural rather than generic, that every component should be determined by the model's own fields, that the construction should have a standard place in gauge theory, and that loop observables should be well defined. Table~\ref{tab:standard} orders the same options by closeness to the value path of standard multi-head attention.

\begin{table}[h]
\centering
\small
\begin{tabular}{>{\raggedright}p{0.12\linewidth}>{\raggedright}p{0.38\linewidth}>{\raggedright\arraybackslash}p{0.40\linewidth}}
\toprule
Rank & Option & Reason \\
\midrule
1 & A, exponential link $U_ie^{X_{ij}}U_j^{-1}$ & Group-valued for every algebra element. The non-communicating part is fixed to the flat transport, so nothing is hidden or unidentifiable. The matrix link is the first-order term, the transported and projected messages are its zeroth and first orders, and loops are defined through inverses. It mirrors the lattice expansion $U=e^{aA}$ \cite{wilson1974,kogut1979}. \\
2 & B, compression of a primitive link $\Pi_i\Theta_{ij}\Pi_j$ & The standard gauge-theoretic picture: links primary, channels as a reduction of structure group, and an exact Levi factorization. It needs a model for $G_{ij}$, and its hidden block is invisible to the message unless the receiver energy reads it. \\
3 & C, head-sum link $U_iOB_{ij}TU_j^{-1}$ & Already computed by the code, and its loops have a clean reading. Group-valuedness is generic and degrades as routing concentrates, and nothing in the construction guarantees invertibility. \\
\bottomrule
\end{tabular}
\caption{Most to least principled.}
\label{tab:principled}
\end{table}

\begin{table}[h]
\centering
\small
\begin{tabular}{>{\raggedright}p{0.12\linewidth}>{\raggedright}p{0.38\linewidth}>{\raggedright\arraybackslash}p{0.40\linewidth}}
\toprule
Rank & Option & Relation to standard attention \\
\midrule
1 & C, head-sum link & Identical to the multi-head value path $\sum_jW_OB_{ij}W_Vx_j$ when the frames are trivial. No new parameters. \\
2 & A, exponential link & Standard attention at first order, plus a flat-transport value term $\sum_j\beta_{ij}\Omega_{ij}\mu_j$ that standard attention lacks, plus higher-order cross-head terms through $TO$. \\
3 & B, compression of a primitive link & Requires a new link field $G_{ij}$. Its nearest standard relative, rotary position encoding, acts on scores rather than on values \cite{su2021roformer}. \\
\bottomrule
\end{tabular}
\caption{Closest to least close to standard attention.}
\label{tab:standard}
\end{table}

The two orders are close to reversed, and Option A sits in the middle of the second while leading the first. Read as a direction for building $\Theta_{ij}$, this suggests Option A as the principled extension and Option C as the model's present content written in link form. Proposition~\ref{prop:charts} qualifies the first half. With a constant generator, Option A is not an extension of the unrouted operator but the same operator in logarithmic coordinates, so Option A adds content only through edge-dependent generators, through chart-natural optimization, or through the boundary and component structure of the group that the affine chart ignores.

\section{Open questions}
\label{sec:open}

Section~\ref{sec:aggregation} settles two questions that the first draft of this note left open. Aggregation under Option A attaches the routing to per-head source labels outside the exponent, and the communication principle makes the receiver the carrier of the complement, so the flat term is a residual on the complement only and vanishes when the heads tile the fiber. What remains open is whether a projected residual stream with $Hd<K$ trains as well as the current additive residual, and whether the self-gate of Eq.~\eqref{eq:self_gate} or a general state-dependent self-link should replace part of the block MLP.

Identifiability under Option B depends on whether the chosen receiver energy reads the hidden block.

Proposition~\ref{prop:charts} turns the chart question into measurements. In a trained model the sign of $\det M$, its smallest singular value, and the number of its negative real eigenvalues, per layer and along the training trajectory, would show whether the learned operators stay in the identity component with a real logarithm, or approach the singular boundary where the affine and logarithmic charts part ways. These quantities need only the trained $O$ and $T$.

The triangle defect of Eq.~\eqref{eq:kl_triangle} is computable in trained models from logged routing costs and normalizers without new machinery, and it would show how close trained routing is to additive along paths.

\section{Numerical checks}
\label{sec:checks}

The script \texttt{link\_origin\_checks/checks.py}, with its output in \texttt{checks\_output.txt}, evaluates each identity in float64 at $K=6$, $H=3$, $d=2$ with random frames and maps. Lemma~\ref{lem:lowrank_exp} holds to a relative error of $5.6\times10^{-16}$, its routed form Eq.~\eqref{eq:exp_routed} to $9.1\times10^{-15}$, the determinant identity to $1.5\times10^{-16}$, and the fixed action on $\ker T$ to $5.6\times10^{-16}$. The split of Eq.~\eqref{eq:split} holds to $7.3\times10^{-16}$, and the first-order defect $(\Theta-\Omega)/s-L$ scales linearly with $s$ ($2.1\times10^{-2}$ at $s=10^{-2}$, $2.2\times10^{-3}$ at $s=10^{-3}$). The block factorization holds to $2.6\times10^{-11}$, limited by the conditioning of the random block $A$, and its determinant identity to $4.3\times10^{-15}$. The loop formulas Eq.~\eqref{eq:headsum_loop} and Eq.~\eqref{eq:exp_loop} hold to $2.7\times10^{-16}$ and $7.4\times10^{-16}$, and the loop trace is unchanged under a frame change. The script \texttt{checks\_aggregation.py}, with output in \texttt{checks\_aggregation\_output.txt}, checks Section~\ref{sec:aggregation}. The complement split of Lemma~\ref{lem:complement} holds to $1.6\times10^{-15}$ at $K=7$ with a rank-4 channel, and $I-P$ fixes $\ker T$ and annihilates $\im O$ to $1.2\times10^{-15}$. Exact enumeration over the head-label tuple reproduces the closed form of Eq.~\eqref{eq:aggregate} with a receiver-carried complement to $4.2\times10^{-16}$. With tiled heads $P=I$ to $6.1\times10^{-16}$, and the Option A message equals the head sum with receiving maps $O\,E_{:,h}$ exactly. The fusion residual of Eq.~\eqref{eq:fusion_residual} holds to $5.6\times10^{-16}$, and the SiLU equals the exponential of its log-sigmoid generator to rounding. The script \texttt{checks\_chart.py}, with output in \texttt{checks\_chart\_output.txt}, checks Section~\ref{sec:charts} at $K=7$, $d=3$ and, for the tiled case, $K=6$. Item 4 of Proposition~\ref{prop:charts} holds to $5.2\times10^{-16}$ and $2.5\times10^{-15}$ in the relative error of Eq.~\eqref{eq:conversion} on the two random trials that satisfy the determinant condition, with $TO=Z$ to $4.4\times10^{-15}$. The embedding $\iota$ is a homomorphism to $2.1\times10^{-16}$, $\exp(OT)=\iota(e^{TO})$ and $I+OCT=\iota(I+C\,TO)$ hold to $5.3\times10^{-16}$ and $3.1\times10^{-16}$, and $\det\iota(B)=\det B$ to $2.2\times10^{-16}$. In the tiled case $O'=e^{X}T^{-1}$ gives $O'T=e^{X}$ to $1.8\times10^{-15}$, and the identity initialization has $\log I=0$. At $X=I$ the exponential and affine charts differ by $36\%$ in relative size. These checks ground the algebra at one small size; they say nothing about training behavior.

\begin{thebibliography}{99}
\bibitem{dennis2026projpred} R.~C.~Dennis, \emph{The projected predictive route: theory and computational path}, technical note, manuscripts/Attention/projected\_predictive.tex, 2026.
\bibitem{dennis2026projgrad} R.~C.~Dennis, \emph{The projected gradient route: receiver energies, entropy, and stable updates}, technical note, manuscripts/Attention/projected\_gradient.tex, 2026.
\bibitem{dennis2026glkv2} R.~C.~Dennis, \emph{Gauge-covariant Gaussian attention through projected matrix links}, manuscript, manuscripts/Attention/GL(K)\_attention\_v2.tex, 2026.
\bibitem{glkreview2026} Referee panel report on the GL(K) attention v2 manuscript, manuscripts/reviews/2026-09-23-glk-v2/REVIEW.md, 2026.
\bibitem{culver1966} W.~J.~Culver, On the existence and uniqueness of the real logarithm of a matrix, \emph{Proceedings of the American Mathematical Society} 17, 1146--1151, 1966.
\bibitem{dennis2026template} R.~C.~Dennis, \emph{Template routing}, technical note, manuscripts/Attention/template\_routing.tex, 2026.
\bibitem{elfwing2018silu} S.~Elfwing, E.~Uchibe, and K.~Doya, Sigmoid-weighted linear units for neural network function approximation in reinforcement learning, \emph{Neural Networks} 107, 3--11, 2018.
\bibitem{he2016resnet} K.~He, X.~Zhang, S.~Ren, and J.~Sun, Deep residual learning for image recognition, \emph{Proceedings of the IEEE Conference on Computer Vision and Pattern Recognition}, 770--778, 2016.
\bibitem{hendrycks2016gelu} D.~Hendrycks and K.~Gimpel, Gaussian error linear units (GELUs), arXiv:1606.08415, 2016.
\bibitem{humphreys1975} J.~E.~Humphreys, \emph{Linear Algebraic Groups}, Graduate Texts in Mathematics 21, Springer, 1975.
\bibitem{kogut1979} J.~B.~Kogut, An introduction to lattice gauge theory and spin systems, \emph{Reviews of Modern Physics} 51, 659--713, 1979.
\bibitem{su2021roformer} J.~Su, Y.~Lu, S.~Pan, A.~Murtadha, B.~Wen, and Y.~Liu, RoFormer: enhanced transformer with rotary position embedding, arXiv:2104.09864, 2021.
\bibitem{vaswani2017} A.~Vaswani, N.~Shazeer, N.~Parmar, J.~Uszkoreit, L.~Jones, A.~N.~Gomez, {\L}.~Kaiser, and I.~Polosukhin, Attention is all you need, \emph{Advances in Neural Information Processing Systems} 30, 2017.
\bibitem{wilson1974} K.~G.~Wilson, Confinement of quarks, \emph{Physical Review D} 10, 2445--2459, 1974.
\bibitem{zhang2005schur} F.~Zhang (ed.), \emph{The Schur Complement and Its Applications}, Springer, 2005.
\end{thebibliography}

\end{document}
